\chapter{LITERATURE REVIEW}
The shift toward high-fidelity occupancy sensing is driven by the failure of traditional hardware to detect stationary presence. Ahmad et al. \cite{Ahmad2021} discusses the integration of hybrid PIR and Radar systems to improve detection reliability in smart buildings. The study finds that while hybrid sensors reduce false triggers, they still exhibit nearly 0.0\% reliability when occupants remain perfectly still. However, the paper lacks a solution for low-power visual confirmation, leaving a gap that only computer vision can fill by analyzing physiological micro-movements.

To deploy vision models on edge hardware, extreme model compression is required. Nagel et al. \cite{nagel2021white} provides a comprehensive white paper on neural network quantization, specifically focusing on 8-bit integer (INT8) Post-Training Quantization. Their research finds that converting models to INT8 reduces memory footprints four-fold with less than a 1\% drop in accuracy. Despite these gains, the paper lacks specific implementation strategies for the 512 KB SRAM limits of the ESP32-S3 \cite{esp32s3_datasheet}, which requires specific instruction-set optimizations.

Further optimization of model response is explored by Wu et al. \cite{wu2016qcnn} through the Q-CNN framework. The paper talks about prioritizing response fidelity matching the output of quantized layers to the original floating-point layers rather than just minimizing weight error. They find that this approach allows for up to 20x model compression while maintaining high classification performance. However, the study lacks a distributed deployment strategy, meaning the model still struggles to process high-resolution images on a single low-power node.

Architectural efficiency is a core theme in reducing parameter counts. Iandola et al. \cite{iandola2016squeezenet} introduces the SqueezeNet architecture, which utilizes Fire Modules to replace expensive filters with $1 \times 1$ squeeze layers. The study finds that SqueezeNet achieves AlexNet-level accuracy with 50x fewer parameters, making it ideal for mobile deployment. While efficient, the architecture lacks the ability to handle extremely low-bit operations, which are necessary for the ultra-low-power requirements of this specific project.

To push efficiency further, Juefei-Xu et al. \cite{juefei2017lbcnn} proposes Local Binary Convolutional Neural Networks (LBCNN). The paper talks about replacing standard learnable filters with fixed, sparse binary kernels that do not require weight updates during training. They find that this method achieves a massive 169x parameter reduction, significantly lowering inference latency. The primary drawback is that LBCNN lacks a native method for processing ultra-high-resolution (4K) frames, which typically causes memory overflows on microcontrollers.

Processing high-resolution data requires spatial decomposition. Li et al. \cite{li2020pytorch} discusses data partitioning strategies to handle large tensors in deep learning workflows. Their findings suggest that splitting images into a grid (tiling) preserves high-frequency spatial details that are usually lost when an image is resized to fit a standard model. However, the paper focuses on high-end GPU clusters and lacks a lightweight communication protocol designed for resource-constrained micro-clusters.

Small object detection is specifically addressed by Akyon et al. \cite{akyon2022sahi} through the SAHI (Slicing Aided Hyper Inference) framework. The paper talks about slicing high-resolution images into overlapping patches to help the model see smaller details at a distance. They find that SAHI increases detection accuracy by 20--30\% for small objects. The main limitation is the high computational overhead, as the framework is designed for powerful workstations rather than real-time edge devices.

Finally, the communication between edge nodes is critical for a distributed system. Jeong et al. \cite{jeong2021} evaluates low-latency communication protocols for IoT device clusters. Their research finds that peer-to-peer protocols can achieve latencies as low as 1.115 ms, which is sufficient for real-time synchronization. However, the study lacks an application-specific test for image tile transmission, which involves higher data throughput than simple sensor packets. This project fills these combined gaps by deploying a real-time, distributed tiling system on an ESP32 cluster.

